\chapter{Технологическая часть}

\section{Средства разработки}

В качестве языка программирования был выбран python3 \cite{python3}, так как в его стандартной библиотеке присутствуют функции замера процессорного времени и использованной памяти, которые требуются в условиях, а также данный язык обладающий множеством инструментов для визуализации и работы с данными и таблицами.

Для построения графиков использовалась библиотека plotly \cite{python3-plotly}.

\section{Реализация алгоритмов}

Ниже представлены реализации алгоритмов линейного поиска и бинарного поиска. Оба алгоритма включают счётчик сравнений в теле цикла: в линейном поиске он увеличивается на 1 за каждую итерацию, в то время как в бинарном поиске счётчик может увеличиваться на 1 или 2, в зависимости от сравнения текущего элемента с искомым значением.

В листинге \ref{lst:lin} представлена реализация алгоритма линейного поиска.

\begin{lstlisting}[label=lst:lin,caption={Алгоритм линейного поиска}]
def FindElem(arr, el):
    cnt = 0
    for i in range(len(arr)):
        cnt += 1
        if arr[i] == el:
            return i, cnt
    return -1, cnt
\end{lstlisting}

\clearpage

В листинге \ref{lst:bin} представлена реализация алгоритма бинарного поиска.

\begin{lstlisting}[label=lst:bin,caption={Алгоритм бинарного поиска}]
def BinFindElem(arr, el):
    left = 0
    right = len(arr) - 1
    cnt = 0

    while left <= right:
        mid = (left + right) // 2
        
        if arr[mid] < el:
            cnt += 1
            left = mid + 1
        elif arr[mid] > el:
            right = mid - 1
            cnt += 2
        else:
            cnt += 2
            return mid, cnt

    return -1, cnt
\end{lstlisting}

\section{Вывод}

В ходе работы были разработаны алгоритмы линейного и бинарного поиска.

\clearpage
